\chapter{How to use this guide}

\method{} is an object detector that looks through two cameras at once, an ordinary visible camera and a
thermal camera, and decides at every location which one to trust. Understanding it end to end touches five
areas that are usually taught separately: deep learning for images, object detection, state-space models,
multispectral fusion and multi-object tracking. This guide puts them in the order the project needs them and,
for each one, names the lectures to watch, the papers to read, the files in the repository that use the idea,
and a short set of exercises.

The path assumes comfortable Python and first-year linear algebra and calculus. It does not assume prior deep
learning. Part time, it takes about two months; most people can start contributing to the code after the
second stage.

\section{The five stages}

\begin{tabularx}{\textwidth}{@{}P{0.6cm}P{4.2cm}P{1.8cm}L@{}}
\hd{} & \hd{stage} & \hd{time} & \hd{what it unlocks in \method{}} \\
1 & Foundations & 3--4 weeks & training loops, convolutions, residual networks, attention \\
2 & Detection and YOLO & 2 weeks & the YOLO26 backbone, neck and NMS-free head the method is built on \\
3 & State-space models and Mamba & 2 weeks & the selective scan and the reliability gate on its step size \\
4 & Seeing in two spectra & 1 week & paired visible--thermal data, fusion designs, the datasets \\
5 & Tracking and video & 1 week & identities over time, the showcase tracker, the next phase \\
\end{tabularx}

\section{How each chapter is organised}

Every stage chapter has the same six parts. \emph{Why it matters} says what breaks in the project without it.
\emph{What to learn} lists the ideas to own, not just recognise. \emph{Lectures and reading} gives the
resources, with links. \emph{Papers} lists the PDFs in \file{study-prep/papers}, each with the one thing to take
from it. \emph{In this repository} names the code to read once the stage is done. \emph{Exercises} are small
and concrete; the chapter ends with a test of readiness.

\section{How to study}

Watch a lecture once at normal speed, then work the matching assignment or exercise before moving on; the
assignments teach more than the videos. Read a paper three times: first the abstract, figures and conclusion;
then the method, with a pen, rederiving each equation; then the experiments, asking what was compared, against
what, and whether the comparison is fair. Keep a notebook of every idea in your own words. When a paper and a
lecture disagree, the code decides.
